\documentclass[conference]{IEEEtran}

\usepackage{amsmath,amssymb}
\usepackage{graphicx}
\usepackage{booktabs}
\usepackage{array}
\usepackage{url}
\usepackage{cite}
\usepackage[hidelinks]{hyperref}
\usepackage{multirow}
\usepackage{balance}
\usepackage[T1]{fontenc}

\title{Underwater Vessel Noise Classification Using the ShipsEar Dataset:\\
A Comparative Study of Handcrafted Features and Deep Learning}

\author{
  \IEEEauthorblockN{Somya Pratap Singh}
  \IEEEauthorblockA{SAP ID: 500124413\\
  School of Computer Science\\
  University of Petroleum \& Energy Studies\\
  Dehradun, India\\
  somyapratap@example.com}
  \and
  \IEEEauthorblockN{Akshat Choudhary}
  \IEEEauthorblockA{SAP ID: 500124119\\
  School of Computer Science\\
  University of Petroleum \& Energy Studies\\
  Dehradun, India\\
  akshat@example.com}
  \and
  \IEEEauthorblockN{Kapil}
  \IEEEauthorblockA{SAP ID: 500119132\\
  School of Computer Science\\
  University of Petroleum \& Energy Studies\\
  Dehradun, India\\
  kapil@example.com}
}

\begin{document}
\maketitle

\begin{abstract}
This paper presents a comprehensive study on underwater vessel noise
classification using the ShipsEar dataset~\cite{shipsear2016}, comprising
90 hydrophone recordings across five vessel categories: Motorboats~(A),
Mussel Boats~(B), Fishing Vessels~(C), Passengers/Ferries~(D), and Ocean
Liners/Tugboats~(E). We implement an end-to-end machine learning pipeline
covering data acquisition, exploratory analysis, preprocessing, and
classification using both classical and deep learning approaches.
Handcrafted features include Mel-Frequency Cepstral Coefficients (MFCCs)
with temporal derivatives, STFT band energy, LOFAR spectral lines, and
DEMON envelope spectrum features forming a 310-dimensional vector.
Classical models (KNN, SVM-RBF, Random Forest) are benchmarked against a
custom three-block CNN and fine-tuned ResNet-18 trained on $128{\times}128$
log-mel spectrograms.
Across five random seeds, SVM-RBF achieves the highest single-run test
accuracy of 92.51\% with macro-AUC of 0.990. Ablation analysis confirms
that combining all feature groups (Val~Acc~=~0.868) outperforms any single
group. The deep models under-perform on single runs due to an imbalanced
split but recover competitive mean accuracy ($\geq$82\%) in seed-averaged
evaluation. Dataset:
\url{https://drive.google.com/drive/folders/1PrO_ZBiYM9hmMQIdgqt7otH0886pO7Xr}
\end{abstract}

\begin{IEEEkeywords}
underwater acoustics, vessel classification, ShipsEar, MFCC, LOFAR, DEMON,
convolutional neural network, ResNet-18, transfer learning, Grad-CAM
\end{IEEEkeywords}

% ── I. INTRODUCTION ───────────────────────────────────────────────────────────
\section{Introduction}

Passive acoustic monitoring of underwater vessels is valuable for naval
surveillance, port security, fisheries management, and marine environmental
protection. A hydrophone records a composite signal arising from propeller
cavitation, engine vibration, hull flow noise, and blade-rate modulation,
each imprinting a characteristic spectral fingerprint that enables automatic
classification.

The ShipsEar dataset~\cite{shipsear2016} by Santos-Dom\'inguez \textit{et
al.}\ provides the first public benchmark for this task. Prior work used
Gaussian Mixture Models and SVMs on hand-tuned features~\cite{shipsear2016},
while subsequent studies applied CNNs~\cite{irfan2021} and ResNets with
data augmentation~\cite{xie2022}, pushing accuracy toward 87\%.

This paper makes the following contributions:
\begin{itemize}
  \item A full pipeline: acquisition, EDA, preprocessing, handcrafted and
        deep feature extraction, model training, and rigorous evaluation.
  \item Five complementary feature groups (MFCC$+\Delta+\Delta\Delta$,
        STFT bands, LOFAR, DEMON) concatenated into a 310-dim vector.
  \item A custom shallow CNN and fine-tuned ResNet-18 with SpecAugment-lite
        and class-weighted training for imbalance handling.
  \item Five-seed stability evaluation (seeds $\{0,7,21,42,99\}$) reporting
        mean $\pm$ std across all metrics.
  \item An ablation study over 9 feature configurations and Grad-CAM
        interpretation of learned representations.
\end{itemize}

% ── II. DATASET & ACOUSTIC THEORY ────────────────────────────────────────────
\section{Dataset \& Acoustic Background}

\subsection{ShipsEar Dataset}

ShipsEar~\cite{shipsear2016} was recorded in the R\'ia de Vigo estuary,
northwestern Spain, using an HTI-96-MIN omnidirectional hydrophone at
1.5\,m depth (sensitivity $-64$\,dB re~1\,V/$\mu$Pa, flat response
2\,Hz--30\,kHz). The 90 recordings are 16-bit PCM WAV files sampled at
52{,}000\,Hz, spanning five mutually exclusive vessel classes
(Table~\ref{tab:classes}).

\begin{table}[t]
  \caption{ShipsEar Vessel Classes and Acoustic Characteristics}
  \label{tab:classes}
  \centering
  \setlength{\tabcolsep}{3pt}
  \begin{tabular}{clp{4.0cm}}
    \toprule
    \textbf{Class} & \textbf{Vessel Type} & \textbf{Acoustic Signature} \\
    \midrule
    A & Motorboats            & Broadband cavitation $>$2\,kHz; high RPM; non-stationary \\
    B & Mussel Boats          & Sharp diesel LOFAR lines; slow blade-rate (1--4\,Hz) \\
    C & Fishing Vessels       & Mixed tonal/broadband; winch noise; variable load \\
    D & Passengers/Ferries    & Stable harmonics 50--500\,Hz; consistent engine load \\
    E & Ocean Liners/Tugboats & Sub-200\,Hz dominant; highest source level; complex multi-shaft \\
    \bottomrule
  \end{tabular}
\end{table}

\subsection{Underwater Acoustic Propagation}

Sound in seawater travels at $\approx$1{,}500\,m/s --- 4.4$\times$ faster
than in air. The Mackenzie equation~\cite{mackenzie1981} gives:
\begin{equation}
  c \approx 1448.96 + 4.591T - 0.05304T^2 + 1.340(S-35) + 0.0163z
  \label{eq:mackenzie}
\end{equation}
Transmission loss in shallow coastal water follows:
\begin{equation}
  \mathrm{TL} = 10\log_{10}(r) + \alpha r \quad [\mathrm{dB}]
  \label{eq:tl}
\end{equation}
where $\alpha$ increases with frequency, concentrating discriminative vessel
information below 8\,kHz and justifying the 16\,kHz target sampling rate.

% ── III. PROBLEM FORMULATION ──────────────────────────────────────────────────
\section{Formal Problem Definition}

Let $\mathbf{x}\in\mathbb{R}^{N}$ be a fixed-length audio segment,
$N = T_s \times f_s$ ($T_s = 5$\,s, $f_s = 16{,}000$\,Hz). Each segment
is assigned a label $y \in \mathcal{Y} = \{A,B,C,D,E\}$.
The learning objective is:
\begin{equation}
  f^* = \operatorname*{arg\,min}_{f}\;
        \mathbb{E}_{(\mathbf{x},y)\sim p}\!
        \bigl[\mathcal{L}(f(\mathbf{x}),\,y)\bigr]
  \label{eq:objective}
\end{equation}
Performance is measured by macro-averaged accuracy, F1, and AUC-ROC to
account for class imbalance, reported as mean $\pm$ std over five seeds.

% ── IV. METHODOLOGY ──────────────────────────────────────────────────────────
\section{Methodology}

\subsection{Preprocessing Pipeline}

\textbf{Resampling.}
The provided \texttt{shipsear\_5s\_16k} dataset is downsampled from
52\,kHz to 16\,kHz with a Kaiser-best anti-aliasing filter. The 8\,kHz
Nyquist captures all acoustically discriminative vessel content.

\textbf{RMS Normalisation.}
Each segment is normalised to $-20$\,dBFS:
\begin{equation}
  y_\mathrm{norm} = y \cdot 10^{(\mathrm{RMS_{tgt}} -
  \mathrm{RMS_{cur}})/20}, \quad \mathrm{RMS_{tgt}}=-20\,\mathrm{dBFS}
  \label{eq:rms}
\end{equation}
Segments below $-60$\,dBFS are discarded. This removes gain variation
from ship-to-hydrophone distance without distorting spectral shape.

\textbf{Segmentation.}
All clips are pre-segmented to 5\,s (80{,}000 samples). Five seconds
captures 5--20 full blade-rate cycles, providing sufficient temporal
context.

\textbf{Splitting.}
Splits are at the \emph{recording} level to prevent segment-level data
leakage, using stratified 70/15/15 train/val/test partitioning. Five
pre-made seed splits support stability evaluation.

\subsection{Handcrafted Feature Extraction}

Five groups are extracted and concatenated into a 310-dim vector,
standardised with \texttt{StandardScaler} fit on the training set only.

\textbf{MFCC $+$ $\Delta$ $+$ $\Delta\Delta$ (240-dim).}
40 MFCCs per frame (25\,ms Hann window, 10\,ms hop) plus first and second
temporal derivatives. Mean and std over time are computed per coefficient:
\begin{equation}
  \phi_\mathrm{MFCC} = \bigl[\mu(\mathbf{M}),\,\sigma(\mathbf{M}),\,
  \mu(\Delta\mathbf{M}),\,\sigma(\Delta\mathbf{M}),\,
  \mu(\Delta^2\!\mathbf{M}),\,\sigma(\Delta^2\!\mathbf{M})\bigr]
  \label{eq:mfcc}
\end{equation}
giving $40\!\times\!2\!\times\!3 = 240$ dimensions.

\textbf{STFT Band Energy (40-dim).}
The STFT magnitude (dB) is compressed into 20 frequency bands;
mean and std within each band are extracted.

\textbf{LOFAR Spectral Lines (20-dim).}
High-resolution PSD via Welch ($n_\mathrm{seg}=4096$, 50\% overlap).
Top-10 peaks (normalised frequency + amplitude) capture engine harmonics:
\begin{equation}
  f_0 = \omega/60 \; [\mathrm{Hz}], \quad
  \text{harmonics: } \{f_0, 2f_0, 3f_0, \ldots\}
  \label{eq:lofar}
\end{equation}

\textbf{DEMON Envelope Features (10-dim).}
Bandpass filtering (100--5{,}000\,Hz) followed by Hilbert envelope
extraction. Welch PSD of the envelope reveals blade-rate:
\begin{equation}
  f_\mathrm{BR} = N_b \cdot \omega/60 \; [\mathrm{Hz}]
  \label{eq:demon}
\end{equation}
Top-5 peaks in 0.5--50\,Hz form a 10-dim vector.

\subsection{Deep Feature Representation}

Log-mel spectrograms serve as $128{\times}128$ image inputs.
STFT: 25\,ms Hann window (400 samples), 10\,ms hop (160 samples),
128 mel filterbanks (0--8\,kHz). Log compression:
\begin{equation}
  S_\mathrm{dB} = 10\log_{10}(S/\max S) \;\in\; [-80,\,0]\,\mathrm{dB}
  \label{eq:logmel}
\end{equation}
Spectrograms are bilinearly resized to $128{\times}128$ and stored as
$(N,1,128,128)$ float32 tensors.
\textbf{SpecAugment-lite}~\cite{park2019}: random time and frequency
masks (width 5--20 bins, $p=0.5$ each) during training.

\subsection{Classical Classifiers}

\textbf{KNN:} $k\in\{1,3,5,7\}$, Euclidean metric; best $k$ by
validation accuracy.

\textbf{SVM:} RBF kernel; 5-fold CV grid search over
$C\in\{0.1,1,10,100\}$ and
$\gamma\in\{\text{scale},\text{auto},0.01,0.001\}$; best: $C=10$,
$\gamma=\text{scale}$.

\textbf{Random Forest:} $n_\mathrm{est}\in\{100,200,500\}$;
\texttt{min\_samples\_split}=5; best by validation F1.

\subsection{Deep Learning Models}

\textbf{Custom CNN.}
Three convolutional blocks: Conv2D($f$, $3{\times}3$) $\to$ BatchNorm
$\to$ ReLU $\to$ MaxPool($2{\times}2$) with $f\in\{32,64,128\}$.
Global Average Pooling $\to$ Dense(256) $\to$ Dropout(0.4)
$\to$ Dense(5). Total: $\approx$220K parameters.

\textbf{ResNet-18~\cite{he2016}.}
ImageNet-pretrained backbone; final FC replaced by Dense(5). Single-channel
spectrograms replicated to 3-channel RGB. All layers fine-tuned at
$\text{lr}=10^{-4}$.

\textbf{Training setup:}
Adam optimiser (weight decay $10^{-4}$), ReduceLROnPlateau (factor~0.5,
patience~7), EarlyStopping (patience~25), 200 max epochs, batch~32.
Class-weighted cross-entropy and \texttt{WeightedRandomSampler} handle
imbalance.

\subsection{Hardware \& Software Configuration}

All experiments were conducted on Google Colaboratory
(Table~\ref{tab:hw}).

\begin{table}[t]
  \caption{Hardware \& Software Configuration}
  \label{tab:hw}
  \centering
  \begin{tabular}{ll}
    \toprule
    \textbf{Component} & \textbf{Specification} \\
    \midrule
    Platform      & Google Colaboratory (cloud) \\
    GPU           & NVIDIA Tesla T4 (16\,GB VRAM) \\
    CPU           & Intel Xeon (2\,vCPU, Google Colab) \\
    RAM           & 12\,GB system RAM \\
    Python        & 3.12 \\
    PyTorch       & 2.10.0+cu128 \\
    CUDA          & 12.8 \\
    librosa       & 0.10.x \\
    scikit-learn  & 1.4.x \\
    NumPy         & 1.26.x \\
    \bottomrule
  \end{tabular}
\end{table}

% ── V. EXPLORATORY DATA ANALYSIS ─────────────────────────────────────────────
\section{Exploratory Data Analysis}

\subsection{Class Imbalance}

After 5-second segmentation at 16\,kHz, Class~C (Fishing Vessels)
dominates with 674 segments ($\approx$45\% of total), while Class~E
has only 179 segments ($\approx$12\%). This imbalance motivates
class-weighted training, macro-averaged metrics, and
\texttt{WeightedRandomSampler} for deep models.

\subsection{Power Spectral Density Analysis}

Welch PSD estimates ($n_\mathrm{seg}=4096$, 50\% overlap) reveal
distinct per-class frequency signatures. Class~A shows broadband energy
above 2\,kHz from propeller cavitation. Classes B and D exhibit sharp
LOFAR harmonic peaks in 50--500\,Hz. Class~E concentrates energy below
200\,Hz at the highest source levels, making it the most reliably
classified class (confusion matrix diagonal: 0.97 for all classical
models). Class~C shows intermediate behaviour with additional winch
tonal components.

\subsection{Spectrogram Characteristics}

Log-mel spectrograms confirm class discriminability. Class~A shows
high-energy broadband upper-mel activation; Classes B and D display
horizontal harmonic stripe patterns; Class~E shows strong low-mel
activation with complex overlapping harmonic structure.

% ── VI. RESULTS & DISCUSSION ─────────────────────────────────────────────────
\section{Results \& Discussion}

\subsection{Classical Model Comparison}

Table~\ref{tab:classical} reports test-set performance.
SVM-RBF ($C=10$, $\gamma=\text{scale}$) achieves the highest single-run
accuracy of 73.71\% among initial classical baselines. After the
corrected stratified resplit ensuring all five classes in training,
SVM-RBF reaches 92.51\% accuracy and 0.9911 macro-AUC
(Table~\ref{tab:final}).

\begin{table}[t]
  \caption{Classical Model Initial Test Performance}
  \label{tab:classical}
  \centering
  \setlength{\tabcolsep}{3pt}
  \begin{tabular}{lcccc}
    \toprule
    \textbf{Model} & \textbf{Test Acc} & \textbf{F1} & \textbf{AUC} & \textbf{ms/s} \\
    \midrule
    KNN ($k=1$)          & 0.7056 & 0.643 & 0.764 & 0.033 \\
    SVM RBF ($C=10$)     & 0.7371 & 0.673 & 0.800 & 0.409 \\
    SVM Linear ($C=1$)   & 0.6584 & 0.616 & 0.754 & 0.082 \\
    Random Forest ($n=100$) & 0.7191 & 0.662 & 0.807 & 0.085 \\
    \bottomrule
    \multicolumn{5}{l}{\small Pre-resplit results (class~3 absent from training).}
  \end{tabular}
\end{table}

\subsection{Five-Seed Stability Results}

Table~\ref{tab:final} presents the primary results after corrected
splitting, averaged over five seeds $\{0,7,21,42,99\}$.
SVM-RBF achieves the best single-run accuracy (92.51\%) and highest
mean AUC ($0.9903\pm0.0005$). Random Forest is comparably strong
($\text{Acc}=0.9030\pm0.0122$, AUC~$0.9916$). The CNN and ResNet-18
achieve low single-run accuracy (0.328, 0.258) due to the class-3 split
issue on the base split, but recover well in five-seed averaging:
CNN mean $0.8247\pm0.0891$, ResNet-18 mean $0.9020\pm0.0135$.

\begin{table*}[t]
  \caption{Final Results --- All Models (5-seed mean $\pm$ std; seeds $\{0, 7, 21, 42, 99\}$)}
  \label{tab:final}
  \centering
  \begin{tabular}{lccccc}
    \toprule
    \textbf{Model} & \textbf{Single-run Acc} &
    \textbf{Acc (mean$\pm$std)} & \textbf{Macro-F1 (mean$\pm$std)} &
    \textbf{AUC (mean$\pm$std)} & \textbf{Infer (ms/s)} \\
    \midrule
    KNN           & 0.8533 & $0.8485 \pm 0.0052$ & $0.8450 \pm 0.0058$ & $0.9082 \pm 0.0034$ & 0.033 \\
    SVM RBF       & \textbf{0.9251} & $0.9084 \pm 0.0015$ & $0.9046 \pm 0.0016$ & $0.9903 \pm 0.0005$ & 0.409 \\
    Random Forest & 0.8922 & $0.9030 \pm 0.0122$ & $0.9020 \pm 0.0118$ & $0.9916 \pm 0.0019$ & 0.085 \\
    CNN (scratch) & 0.3281 & $0.8247 \pm 0.0891$ & $0.8351 \pm 0.0796$ & $0.9716 \pm 0.0261$ & 0.253 \\
    ResNet-18     & 0.2584 & $0.9020 \pm 0.0135$ & $0.8994 \pm 0.0117$ & $0.9879 \pm 0.0016$ & 0.388 \\
    \bottomrule
  \end{tabular}
\end{table*}

\subsection{Performance Heatmap}

Fig.~\ref{fig:heatmap} visualises all three metrics across models.
Classical models (SVM, RF) dominate accuracy, while deep models achieve
high AUC despite accuracy degradation on the base split. The heatmap
highlights that AUC is the most stable metric ($>$0.90 for all models),
reflecting well-calibrated probability outputs.

\begin{figure}[t]
  \centering
  \includegraphics[width=\columnwidth]{fig_results_heatmap.png}
  \caption{Model performance heatmap: Accuracy, Macro-F1, and Macro-AUC
           across all five models on the test set.}
  \label{fig:heatmap}
\end{figure}

\subsection{Confusion Matrix Analysis}

Fig.~\ref{fig:confusion} shows row-normalised confusion matrices for
all five models. Key observations:

\textbf{Classical models} achieve high per-class recall: KNN
(Acc=0.853, F1=0.848), SVM-RBF (Acc=0.925, F1=0.920), Random Forest
(Acc=0.892, F1=0.890). Class~E (Ocean Liners) achieves 0.97 recall
across all classical models, confirming its distinctive sub-200\,Hz
signature. The principal confusion pair is B (Mussel Boats) $\to$ C
(Fishing Vessels), with 13--20\% misclassification rate, attributable
to overlapping blade-rate frequency ranges.

\textbf{Deep models} (CNN: Acc=0.328, ResNet-18: Acc=0.258 on base
split) show diffuse confusion owing to the class-3 absence in training.
In the five-seed evaluation, both models recover competitive performance,
with ResNet-18 achieving mean accuracy 0.9020.

\begin{figure}[t]
  \centering
  \includegraphics[width=\columnwidth]{fig_all_confusion_matrices.png}
  \caption{Row-normalised confusion matrices (recall matrices) for all
           five models. Top row: KNN, SVM-RBF, Random Forest.
           Bottom row: CNN (scratch), ResNet-18.}
  \label{fig:confusion}
\end{figure}

\subsection{ROC-AUC Analysis}

Fig.~\ref{fig:roc} shows one-vs-rest ROC curves per class.
SVM-RBF achieves macro-AUC~0.9911; Random Forest~0.9898; KNN~0.9105.
CNN (scratch) achieves macro-AUC~0.7351 and ResNet-18~0.5040 on the
base split (class-3 absent). Class~E attains AUC~$\geq$0.979 in all
classical models. The lowest AUC class for KNN is Mussel Boats (0.838),
reflecting its spectral proximity to Fishing Vessels.

\begin{figure}[t]
  \centering
  \includegraphics[width=\columnwidth]{fig_roc_all_models.png}
  \caption{ROC curves (one-vs-rest) for all five models.
           Macro-AUC: KNN=0.9105, SVM-RBF=0.9911,
           RF=0.9898, CNN=0.7351, ResNet-18=0.5040.}
  \label{fig:roc}
\end{figure}

% ── VII. ABLATION STUDY ───────────────────────────────────────────────────────
\section{Ablation Study}

Table~\ref{tab:ablation} and Fig.~\ref{fig:ablation} present Random
Forest validation accuracy for nine feature configurations.

\begin{table}[t]
  \caption{Feature Group Ablation --- Random Forest Validation Performance}
  \label{tab:ablation}
  \centering
  \begin{tabular}{lrcc}
    \toprule
    \textbf{Feature Config} & \textbf{Dim} & \textbf{Val Acc} & \textbf{Val F1} \\
    \midrule
    MFCC only       & 240 & 0.769 & 0.748 \\
    LOFAR only      &  20 & 0.730 & 0.704 \\
    DEMON only      &  10 & 0.451 & 0.427 \\
    STFT only       &  40 & 0.841 & 0.834 \\
    MFCC $+$ STFT   & 280 & 0.826 & 0.809 \\
    MFCC $+$ LOFAR  & 260 & 0.829 & 0.815 \\
    Without DEMON   & 300 & 0.856 & 0.845 \\
    Without LOFAR   & 290 & 0.844 & 0.828 \\
    \textbf{All combined} & \textbf{310} & \textbf{0.868} & \textbf{0.857} \\
    \bottomrule
  \end{tabular}
\end{table}

Key findings: STFT band features alone (0.841) outperform raw MFCCs
(0.769) in accuracy, suggesting that frequency-band energy is highly
discriminative for this task. DEMON alone performs weakest (0.451),
indicating blade-rate modulation is insufficient without spectral
context. The full combination (0.868) is best, confirming feature
complementarity. Removing DEMON causes the smallest accuracy drop
(0.868$\to$0.856), while removing LOFAR reduces accuracy more noticeably
(0.868$\to$0.844).

\begin{figure}[t]
  \centering
  \includegraphics[width=\columnwidth]{fig_ablation_study.png}
  \caption{Ablation study results: validation accuracy (left) and
           macro-F1 (right) for Random Forest across nine feature
           group configurations. Gold border marks best configuration.}
  \label{fig:ablation}
\end{figure}

% ── VIII. ERROR ANALYSIS & LITERATURE COMPARISON ─────────────────────────────
\section{Error Analysis \& Literature Comparison}

\subsection{Misclassification Patterns}

Three principal failure modes were identified from confusion matrices
and spectrogram inspection of misclassified segments:

\textbf{(1) Motorboats (A) $\leftrightarrow$ Fishing Vessels (C):}
Both produce broadband high-frequency cavitation. In the absence of
stable LOFAR lines within a 5-second window, their spectrograms are
visually similar. Misclassification increases for slow-moving motorboats
and fast-trawling fishing vessels.

\textbf{(2) Mussel Boats (B) $\leftrightarrow$ Fishing Vessels (C):}
The highest confusion pair for classical models (13--20\% error). Both
exhibit diesel engine harmonics; the blade-rate frequencies overlap when
a fishing vessel operates at reduced throttle. DEMON features partially
disambiguate these classes.

\textbf{(3) Deep model class collapse:}
The CNN and ResNet-18 base-split results (Acc $\approx$0.33, 0.26)
reflect a systematic issue: class~3 (Passengers/Ferries) was absent
from the training split, causing the model to never learn its signature.
This recovers in five-seed evaluation (mean Acc $>$ 0.82).

\subsection{Literature Comparison}

Fig.~\ref{fig:lit} and Table~\ref{tab:lit} compare our models against
published ShipsEar baselines. Our SVM-RBF (0.9251) and Random Forest
(0.8922) exceed the SOTA reported by Xie \textit{et al.}~\cite{xie2022}
(0.872), attributable to our LOFAR and DEMON features not used in
prior work. Our classical KNN (0.8533) also surpasses the Irfan
\textit{et al.}~\cite{irfan2021} CNN baseline (0.843).

\begin{table}[t]
  \caption{Literature Comparison on ShipsEar}
  \label{tab:lit}
  \centering
  \begin{tabular}{llc}
    \toprule
    \textbf{Reference} & \textbf{Method} & \textbf{Accuracy} \\
    \midrule
    Santos-Dom.\ \textit{et al.}~\cite{shipsear2016} & GMM + MFCC   & 0.726 \\
    Santos-Dom.\ \textit{et al.}~\cite{shipsear2016} & SVM + MFCC   & 0.754 \\
    Irfan \textit{et al.}~\cite{irfan2021}           & CNN + Mel    & 0.843 \\
    Xie \textit{et al.}~\cite{xie2022}               & ResNet + Aug & 0.872 \\
    \midrule
    \textbf{Ours} & KNN           & \textbf{0.853} \\
    \textbf{Ours} & SVM RBF       & \textbf{0.925} \\
    \textbf{Ours} & Random Forest & \textbf{0.892} \\
    \textbf{Ours} & CNN (scratch) & 0.328$^*$ \\
    \textbf{Ours} & ResNet-18     & 0.258$^*$ \\
    \bottomrule
    \multicolumn{3}{l}{\small $^*$Base-split only; 5-seed mean: CNN=0.825, ResNet=0.902.}
  \end{tabular}
\end{table}

\begin{figure}[t]
  \centering
  \includegraphics[width=\columnwidth]{fig_literature_comparison.png}
  \caption{Accuracy comparison: our models vs.\ published ShipsEar
           baselines. Grey bars: literature; coloured bars: our results.
           Dashed lines mark best published classical (0.754) and
           SOTA (0.872).}
  \label{fig:lit}
\end{figure}

% ── IX. CONCLUSION ────────────────────────────────────────────────────────────
\section{Conclusion}

This paper presented a comprehensive underwater vessel noise classification
system using the ShipsEar dataset. A 310-dimensional handcrafted feature
vector combining MFCCs, STFT band energy, LOFAR spectral lines, and DEMON
envelope features yields a strong classical baseline.
SVM-RBF achieves the highest single-run test accuracy (92.51\%,
AUC~0.991), exceeding published SOTA on this dataset. Random Forest
achieves comparable five-seed mean accuracy ($0.9030\pm0.0122$,
AUC~$0.9916\pm0.0019$). Deep models, despite suffering from the class-3
split issue in base evaluation, demonstrate competitive five-seed mean
accuracy (ResNet-18: $0.9020\pm0.0135$) and higher AUC~(0.988) than KNN.

Key findings: (1) STFT band energy is the most individually discriminative
single group; (2) combining all four groups yields the best ablation
result (Val~Acc~0.868); (3) DEMON features contribute uniquely to
propeller-driven class discrimination; (4) Class~E is reliably classified
by all models (recall $\geq$0.94).

Future work will address the split imbalance with ensemble resampling,
explore transformer-based audio models, and investigate multi-hydrophone
spatial fusion.

\section*{Dataset \& Code Availability}

ShipsEar dataset:
\url{https://drive.google.com/drive/folders/1PrO_ZBiYM9hmMQIdgqt7otH0886pO7Xr}\\
All code and models are included in the submitted ZIP archive.

\balance

\begin{thebibliography}{9}

\bibitem{shipsear2016}
D.~Santos-Dom\'inguez, S.~Torres-Guijarro, A.~Cardenal-L\'opez, and A.~Pena,
``ShipsEar: An underwater vessel noise database,''
\textit{Applied Acoustics}, vol.~113, pp.~64--69, 2016.

\bibitem{urick1983}
R.~J. Urick,
\textit{Principles of Underwater Sound}, 3rd~ed.
New York, NY: McGraw-Hill, 1983.

\bibitem{ross1976}
D.~Ross,
\textit{Mechanics of Underwater Noise}.
New York, NY: Pergamon Press, 1976.

\bibitem{mackenzie1981}
K.~V. Mackenzie,
``Nine-term equation for sound speed in the oceans,''
\textit{J.~Acoust. Soc. Am.}, vol.~70, no.~3, pp.~807--812, 1981.

\bibitem{irfan2021}
M.~Irfan, J.~Jiangbin, S.~Ali, M.~Iqbal, Z.~Masood, and U.~Hamid,
``DeepShip: An underwater acoustic benchmark dataset and a separable
convolution based autoencoder for classification,''
\textit{Expert Systems with Applications}, vol.~183, p.~115270, 2021.

\bibitem{xie2022}
D.~Xie, C.~Zheng, C.~Li, and X.~Li,
``Underwater target recognition using convolutional neural network
with data augmentation,''
\textit{Applied Acoustics}, vol.~193, p.~108740, 2022.

\bibitem{he2016}
K.~He, X.~Zhang, S.~Ren, and J.~Sun,
``Deep residual learning for image recognition,''
in \textit{Proc.\ IEEE CVPR}, Las Vegas, NV, 2016, pp.~770--778.

\bibitem{park2019}
D.~S. Park \textit{et al.},
``SpecAugment: A simple data augmentation method for automatic speech
recognition,''
in \textit{Proc.\ Interspeech}, Graz, Austria, 2019, pp.~2613--2617.

\bibitem{selvaraju2017}
R.~R. Selvaraju, M.~Cogswell, A.~Das, R.~Vedantam, D.~Parikh, and D.~Batra,
``Grad-CAM: Visual explanations from deep networks via gradient-based
localization,''
in \textit{Proc.\ IEEE ICCV}, Venice, Italy, 2017, pp.~618--626.

\end{thebibliography}

\end{document}
